\color{unidarkgreen}{\section[(обязательное) Уставки функций РЗА]{(обязательное)\\Уставки функций РЗА}\label{app:set}}
\color{black}

\setlength{\extrarowheight}{0.06cm}

\par
Уставки ДЗТ представлены в таблице \ref{app_set:dzt}.

\footnotesize
\begin{longtable}{|>{\centering\arraybackslash}m{5.3cm}|>{\centering\arraybackslash}m{3.3cm}|>{\centering\arraybackslash}m{4.2cm}|>{\centering\arraybackslash}m{1.8cm}|>{\centering\arraybackslash}m{1cm}|}
% Первая страница: полный заголовок
\caption{Параметры для настройки ДЗТ\hfill\vspace{-0.5\baselineskip}}\label{app_set:dzt}\\
\hline
\rowcolor{gray!20}
Параметр & Условное обозначение на схеме & Значение/ Диапазон & Единица измерения & Шаг \\
\hline
\endfirsthead

% Промежуточные страницы: «Продолжение таблицы»
\caption*{\hspace{3pt}\emph{Продолжение таблицы \ref{app_set:dzt}\hfill\vspace{-0.5\baselineskip}}} \\ % сделано по ГОСТ 2.105 п.6.8.7
\hline
\rowcolor{gray!20}
Параметр & Условное обозначение на схеме & Значение/ Диапазон & Единица измерения & Шаг \\
\hline
\endhead

%==+t1*PDIF1|TDIF
\multicolumn{5}{|c|}{ Общие } \\ \hline 
\centering Использование стороны 1 (Сторона1)  & \centering Сторона1 & \centering Не используется \\ Используется & \centering -- & \centering \arraybackslash -- \\
\hline
\centering Использование стороны 2 (Сторона2)  & \centering Сторона2 & \centering Не используется \\ Используется & \centering -- & \centering \arraybackslash -- \\
\hline
\centering Использование стороны 3 (Сторона3)  & \centering Сторона3 & \centering Не используется \\ Используется & \centering -- & \centering \arraybackslash -- \\
\hline
\centering Базисная мощность (Sбаз)  & \centering Sбаз & \centering 1000...500000 & \centering кВА & \centering \arraybackslash 100 \\
\hline
\multicolumn{5}{|c|}{ Сторона 1 } \\ \hline 
\centering Схема соединения трансформаторов тока стороны 1 (КсхемСт1)  & \centering КсхемСт1 & \centering Звезда \\ Треугольник & \centering -- & \centering \arraybackslash -- \\
\hline
\centering Номинальное напряжение стороны 1 (Uном1)  & \centering Uном1 & \centering 6,00...250,00 & \centering кВ & \centering \arraybackslash 0,01 \\
\hline
\centering Номинальный первичный ток ТТ стороны 1 (IпервСт1)  & \centering IпервСт1 & \centering 1...40000 & \centering А & \centering \arraybackslash 1 \\
\hline
\centering Номинальный вторичный ток ТТ стороны 1 (IвторСт1)  & \centering IвторСт1 & \centering 0,2...5,0 & \centering А & \centering \arraybackslash 0,1 \\
\hline
\centering Коэффициент коррекции полярности ТТ стороны 1 (KпСт1)  & \centering KпСт1 & \centering Синфазное включение & \centering -- & \centering \arraybackslash -- \\
\hline
\centering Компенсация токов 3I0 стороны 1 (Компенс3I0Ст1)  & \centering Компенс3I0Ст1 & \centering Без компенсации токов 3I0 \\ С компенсацией токов 3I0 & \centering -- & \centering \arraybackslash -- \\
\hline
\centering Векторная группа стороны 1 (NсхСт1)  & \centering NсхСт1 & \centering 0...11 & \centering -- & \centering \arraybackslash 1 \\
\hline
\multicolumn{5}{|c|}{ Сторона 2 } \\ \hline 
\centering Схема соединения трансформаторов тока стороны 2 (КсхемСт2)  & \centering КсхемСт2 & \centering Звезда \\ Треугольник & \centering -- & \centering \arraybackslash -- \\
\hline
\centering Номинальное напряжение стороны 2 (Uном2)  & \centering Uном2 & \centering 6,00...250,00 & \centering кВ & \centering \arraybackslash 0,01 \\
\hline
\centering Номинальный первичный ток ТТ стороны 2 (IпервСт2)  & \centering IпервСт2 & \centering 1...40000 & \centering А & \centering \arraybackslash 1 \\
\hline
\centering Номинальный вторичный ток ТТ стороны 2 (IвторСт2)  & \centering IвторСт2 & \centering 0,2...5,0 & \centering А & \centering \arraybackslash 0,1 \\
\hline
\centering Коэффициент коррекции полярности ТТ стороны 2 (KпСт2)  & \centering KпСт2 & \centering Синфазное включение & \centering -- & \centering \arraybackslash -- \\
\hline
\centering Компенсация токов 3I0 стороны 2 (Компенс3I0Ст2)  & \centering Компенс3I0Ст2 & \centering Без компенсации токов 3I0 \\ С компенсацией токов 3I0 & \centering -- & \centering \arraybackslash -- \\
\hline
\centering Векторная группа стороны 2 (NсхСт2)  & \centering NсхСт2 & \centering 0...11 & \centering -- & \centering \arraybackslash 1 \\
\hline
\multicolumn{5}{|c|}{ Сторона 3 } \\ \hline 
\centering Схема соединения трансформаторов тока стороны 3 (КсхемСт3)  & \centering КсхемСт3 & \centering Звезда \\ Треугольник & \centering -- & \centering \arraybackslash -- \\
\hline
\centering Номинальное напряжение стороны 3 (Uном3)  & \centering Uном3 & \centering 6,00...250,00 & \centering кВ & \centering \arraybackslash 0,01 \\
\hline
\centering Номинальный первичный ток ТТ стороны 3 (IпервСт3)  & \centering IпервСт3 & \centering 1...40000 & \centering А & \centering \arraybackslash 1 \\
\hline
\centering Номинальный вторичный ток ТТ стороны 3 (IвторСт3)  & \centering IвторСт3 & \centering 0,2...5,0 & \centering А & \centering \arraybackslash 0,1 \\
\hline
\centering Коэффициент коррекции полярности ТТ стороны 3 (KпСт3)  & \centering KпСт3 & \centering Синфазное включение & \centering -- & \centering \arraybackslash -- \\
\hline
\centering Компенсация токов 3I0 стороны 3 (Компенс3I0Ст3)  & \centering Компенс3I0Ст3 & \centering Без компенсации токов 3I0 \\ С компенсацией токов 3I0 & \centering -- & \centering \arraybackslash -- \\
\hline
\centering Векторная группа стороны 3 (NсхСт3)  & \centering NсхСт3 & \centering 0...11 & \centering -- & \centering \arraybackslash 1 \\
\hline
\multicolumn{5}{|c|}{ ДТЗт } \\ \hline 
\centering Ввод функции в работу (Ввод\_функции)  & \centering SGF1 & \centering Не предусмотрено \\ Предусмотрено & \centering -- & \centering \arraybackslash -- \\
\hline
\centering Контроль от БВКЗ (Контр\_БВКЗ)  & \centering SGF3 & \centering Без контроля БВКЗ \\ С контролем БВКЗ & \centering -- & \centering \arraybackslash -- \\
\hline
\centering Режим блокировки (Реж\_блок)  & \centering SGF2 & \centering Без блокировки \\ Блокировка по 2 гармонике \\ Блокировка по 5 гармонике \\ Блокировка по 2 и 5 гармоникам & \centering -- & \centering \arraybackslash -- \\
\hline
\centering Начальный дифференциальный ток срабатывания (Iср)  & \centering -- & \centering 0,20...0,60 & \centering о.е. & \centering \arraybackslash 0,01 \\
\hline
\centering Начальный дифференциальный ток срабатывания в режиме загрубления (при неисправностях в токовых цепях) (Iср\_загр)  & \centering -- & \centering 0,20...10,00 & \centering о.е. & \centering \arraybackslash 0,01 \\
\hline
\centering Коэффициент торможения первого наклонного участка (Kт1)  & \centering -- & \centering 0,20...1,00 & \centering о.е. & \centering \arraybackslash 0,01 \\
\hline
\centering Начальный тормозной ток первого наклонного участка (Iт1)  & \centering -- & \centering 0,60...3,00 & \centering о.е. & \centering \arraybackslash 0,01 \\
\hline
\centering Коэффициент торможения второго наклонного участка (Kт2)  & \centering -- & \centering 0,20...1,00 & \centering о.е. & \centering \arraybackslash 0,01 \\
\hline
\centering Начальный тормозной ток второго наклонного участка (Iт2)  & \centering -- & \centering 1,20...10,00 & \centering о.е. & \centering \arraybackslash 0,01 \\
\hline
\centering Выдержка времени срабатывания (Tср)  & \centering T1 & \centering 0...3000 & \centering мс & \centering \arraybackslash 10 \\
\hline
\multicolumn{5}{|c|}{ ДТО } \\ \hline 
\centering Ввод функции в работу (Ввод\_функции)  & \centering SGF1 & \centering Не предусмотрено \\ Предусмотрено & \centering -- & \centering \arraybackslash -- \\
\hline
\centering Дифференциальный ток срабатывания (Iср)  & \centering I> & \centering 3,0...20,0 & \centering о.е. & \centering \arraybackslash 0,1 \\
\hline
\centering Выдержка времени срабатывания (Tср)  & \centering T1 & \centering 0...3000 & \centering мс & \centering \arraybackslash 10 \\
\hline
\multicolumn{5}{|c|}{ Д2Г } \\ \hline 
\centering Коэффициент отношения второй гармоники к первой гармонике для выполнения пусковых условий (Ih2/Ih1)  & \centering I2г/I1г> & \centering 0,10...0,50 & \centering -- & \centering \arraybackslash 0,01 \\
\hline
\centering Режим перекрестной блокировки (Перекр\_блок)  & \centering SGF1 & \centering Не предусмотрено \\ Предусмотрено & \centering -- & \centering \arraybackslash -- \\
\hline
\centering Выдержка времени ввода перекрестной блокировки (Tперек\_блок)  & \centering T2 & \centering 60...4000 & \centering мс & \centering \arraybackslash 10 \\
\hline
\centering Выдержка времени на возврат (Tвоз)  & \centering T1 & \centering 0...100 & \centering мс & \centering \arraybackslash 1 \\
\hline
\multicolumn{5}{|c|}{ Д5Г } \\ \hline 
\centering Коэффициент отношения пятой гармоники к первой гармонике для выполнения пусковых условий (Ih5/Ih1)  & \centering I5г/I1г> & \centering 0,05...0,40 & \centering -- & \centering \arraybackslash 0,01 \\
\hline
\centering Режим перекрестной блокировки (Перекр\_блок)  & \centering SGF1 & \centering Не предусмотрено \\ Предусмотрено & \centering -- & \centering \arraybackslash -- \\
\hline
\centering Выдержка времени ввода перекрестной блокировки (Tперек\_блок)  & \centering T2 & \centering 60...4000 & \centering мс & \centering \arraybackslash 10 \\
\hline
\centering Выдержка времени на возврат (Tвоз)  & \centering T1 & \centering 0...100 & \centering мс & \centering \arraybackslash 1 \\
\hline
\multicolumn{5}{|c|}{ БВКЗ } \\ \hline 
\centering Работа селектора по отрицательной/положительной полуволне (Знак\_полуволны)  & \centering -- & \centering Отрицательная полуволна \\ Положительная полуволна & \centering -- & \centering \arraybackslash -- \\
\hline
\centering Зона нечувствительности (Зона\_нечувств)  & \centering -- & \centering 0,01...2,00 & \centering о.е. & \centering \arraybackslash 0,01 \\
\hline
\centering Время, соответствующее углу блокировки (TуголБлок)  & \centering -- & \centering 0,0010...0,0090 & \centering с & \centering \arraybackslash 0,0001 \\
\hline
\multicolumn{5}{|c|}{ КЦТнеб } \\ \hline 
\centering Ввод функции в работу (Ввод\_функции)  & \centering SGF1 & \centering Не предусмотрено \\ Предусмотрено & \centering -- & \centering \arraybackslash -- \\
\hline
\centering Ток срабатывания (Iср)  & \centering Idiff> & \centering 0,04...2,00 & \centering о.е. & \centering \arraybackslash 0,01 \\
\hline
\centering Выдержка времени срабатывания (Tср)  & \centering T1 & \centering 0,0...110,0 & \centering с & \centering \arraybackslash 0,1 \\
\hline
%===t1
\end{longtable}
\normalsize

%%%%%%%%%%%%%%%%%%%%%%%%%%%%%%%%%%%%%%%%%%%%%%%%%%%%%%%%%% КЦТ %%%%%%%%%%%%%%%%%%%%%%%%%%%%%%%%%%%%%%%%%%%%%%%%%%%%%%%%%%%%%%%%%%


\par
Уставки КЦТ представлены в таблице \ref{app_set:kzt}.

\footnotesize
\begin{longtable}{|>{\centering\arraybackslash}m{5.3cm}|>{\centering\arraybackslash}m{3.3cm}|>{\centering\arraybackslash}m{4.2cm}|>{\centering\arraybackslash}m{1.8cm}|>{\centering\arraybackslash}m{1cm}|}

\caption{Параметры для настройки КЦТ\hfill\vspace{-0.5\baselineskip}}\label{app_set:kzt}\\
\hline
\rowcolor{gray!20}
Параметр & Условное обозначение на схеме & Значение/ Диапазон & Единица измерения & Шаг \\
\hline
\endfirsthead

\caption*{\hspace{3pt}\emph{Продолжение таблицы \ref{app_set:kzt}\hfill\vspace{-0.5\baselineskip}}} \\ % сделано по ГОСТ 2.105 п.6.8.7
\hline
\rowcolor{gray!20}
Параметр & Условное обозначение на схеме & Значение/ Диапазон & Единица измерения & Шаг \\
\hline
\endhead

%==+t1*RCTR1|CTR_UIRZ
\multicolumn{5}{|c|}{ КЦТст1 } \\ \hline 
\centering Ввод функции в работу (Ввод\_функции)  & \centering SGF1 & \centering Не предусмотрено \\ Предусмотрено & \centering -- & \centering \arraybackslash -- \\
\hline
\centering Режим контроля обрыва провода (Контр\_обр\_пров)  & \centering SGF2 & \centering Не предусмотрено \\ Предусмотрено & \centering -- & \centering \arraybackslash -- \\
\hline
\centering Номинальный ток токового входа терминала (Iном)  & \centering -- & \centering 0,2...5,0 & \centering А & \centering \arraybackslash 0,1 \\
\hline
\centering Минимальное значение фазного тока (Iмин)  & \centering Iмин< & \centering 0,05...1,00 & \centering о.е. & \centering \arraybackslash 0,01 \\
\hline
\centering Коэффициент симметрии (Kсим)  & \centering Iмин/Iмакс< & \centering 0,10...0,95 & \centering о.е. & \centering \arraybackslash 0,01 \\
\hline
\centering Минимальная величина максимального из фазных токов (LIсим)  & \centering Iмакс> & \centering 0,01...2,00 & \centering о.е. & \centering \arraybackslash 0,01 \\
\hline
\centering Выдержка времени на срабатывание по критерию обрыва фазы (Tобрыв)  & \centering T1 & \centering 0,00...100,00 & \centering с & \centering \arraybackslash 0,01 \\
\hline
\centering Выдержка времени на срабатывание по асимметрии токов (Tасим)  & \centering T2 & \centering 0,00...100,00 & \centering с & \centering \arraybackslash 0,01 \\
\hline
\multicolumn{5}{|c|}{ КЦТст2 } \\ \hline 
\centering Ввод функции в работу (Ввод\_функции)  & \centering SGF1 & \centering Не предусмотрено \\ Предусмотрено & \centering -- & \centering \arraybackslash -- \\
\hline
\centering Режим контроля обрыва провода (Контр\_обр\_пров)  & \centering SGF2 & \centering Не предусмотрено \\ Предусмотрено & \centering -- & \centering \arraybackslash -- \\
\hline
\centering Номинальный ток токового входа терминала (Iном)  & \centering -- & \centering 0,2...5,0 & \centering А & \centering \arraybackslash 0,1 \\
\hline
\centering Минимальное значение фазного тока (Iмин)  & \centering Iмин< & \centering 0,05...1,00 & \centering о.е. & \centering \arraybackslash 0,01 \\
\hline
\centering Коэффициент симметрии (Kсим)  & \centering Iмин/Iмакс< & \centering 0,10...0,95 & \centering о.е. & \centering \arraybackslash 0,01 \\
\hline
\centering Минимальная величина максимального из фазных токов (LIсим)  & \centering Iмакс> & \centering 0,01...2,00 & \centering о.е. & \centering \arraybackslash 0,01 \\
\hline
\centering Выдержка времени на срабатывание по критерию обрыва фазы (Tобрыв)  & \centering T1 & \centering 0,00...100,00 & \centering с & \centering \arraybackslash 0,01 \\
\hline
\centering Выдержка времени на срабатывание по асимметрии токов (Tасим)  & \centering T2 & \centering 0,00...100,00 & \centering с & \centering \arraybackslash 0,01 \\
\hline
\multicolumn{5}{|c|}{ КЦТст3 } \\ \hline 
\centering Ввод функции в работу (Ввод\_функции)  & \centering SGF1 & \centering Не предусмотрено \\ Предусмотрено & \centering -- & \centering \arraybackslash -- \\
\hline
\centering Режим контроля обрыва провода (Контр\_обр\_пров)  & \centering SGF2 & \centering Не предусмотрено \\ Предусмотрено & \centering -- & \centering \arraybackslash -- \\
\hline
\centering Номинальный ток токового входа терминала (Iном)  & \centering -- & \centering 0,2...5,0 & \centering А & \centering \arraybackslash 0,1 \\
\hline
\centering Минимальное значение фазного тока (Iмин)  & \centering Iмин< & \centering 0,05...1,00 & \centering о.е. & \centering \arraybackslash 0,01 \\
\hline
\centering Коэффициент симметрии (Kсим)  & \centering Iмин/Iмакс< & \centering 0,10...0,95 & \centering о.е. & \centering \arraybackslash 0,01 \\
\hline
\centering Минимальная величина максимального из фазных токов (LIсим)  & \centering Iмакс> & \centering 0,01...2,00 & \centering о.е. & \centering \arraybackslash 0,01 \\
\hline
\centering Выдержка времени на срабатывание по критерию обрыва фазы (Tобрыв)  & \centering T1 & \centering 0,00...100,00 & \centering с & \centering \arraybackslash 0,01 \\
\hline
\centering Выдержка времени на срабатывание по асимметрии токов (Tасим)  & \centering T2 & \centering 0,00...100,00 & \centering с & \centering \arraybackslash 0,01 \\
\hline
\multicolumn{5}{|c|}{ КЦТст4 } \\ \hline 
\centering Ввод функции в работу (Ввод\_функции)  & \centering SGF1 & \centering Не предусмотрено \\ Предусмотрено & \centering -- & \centering \arraybackslash -- \\
\hline
\centering Режим контроля обрыва провода (Контр\_обр\_пров)  & \centering SGF2 & \centering Не предусмотрено \\ Предусмотрено & \centering -- & \centering \arraybackslash -- \\
\hline
\centering Номинальный ток токового входа терминала (Iном)  & \centering -- & \centering 0,2...5,0 & \centering А & \centering \arraybackslash 0,1 \\
\hline
\centering Минимальное значение фазного тока (Iмин)  & \centering Iмин< & \centering 0,05...1,00 & \centering о.е. & \centering \arraybackslash 0,01 \\
\hline
\centering Коэффициент симметрии (Kсим)  & \centering Iмин/Iмакс< & \centering 0,10...0,95 & \centering о.е. & \centering \arraybackslash 0,01 \\
\hline
\centering Минимальная величина максимального из фазных токов (LIсим)  & \centering Iмакс> & \centering 0,01...2,00 & \centering о.е. & \centering \arraybackslash 0,01 \\
\hline
\centering Выдержка времени на срабатывание по критерию обрыва фазы (Tобрыв)  & \centering T1 & \centering 0,00...100,00 & \centering с & \centering \arraybackslash 0,01 \\
\hline
\centering Выдержка времени на срабатывание по асимметрии токов (Tасим)  & \centering T2 & \centering 0,00...100,00 & \centering с & \centering \arraybackslash 0,01 \\
\hline
\multicolumn{5}{|c|}{ КЦТст5 } \\ \hline 
\centering Ввод функции в работу (Ввод\_функции)  & \centering SGF1 & \centering Не предусмотрено \\ Предусмотрено & \centering -- & \centering \arraybackslash -- \\
\hline
\centering Режим контроля обрыва провода (Контр\_обр\_пров)  & \centering SGF2 & \centering Не предусмотрено \\ Предусмотрено & \centering -- & \centering \arraybackslash -- \\
\hline
\centering Номинальный ток токового входа терминала (Iном)  & \centering -- & \centering 0,2...5,0 & \centering А & \centering \arraybackslash 0,1 \\
\hline
\centering Минимальное значение фазного тока (Iмин)  & \centering Iмин< & \centering 0,05...1,00 & \centering о.е. & \centering \arraybackslash 0,01 \\
\hline
\centering Коэффициент симметрии (Kсим)  & \centering Iмин/Iмакс< & \centering 0,10...0,95 & \centering о.е. & \centering \arraybackslash 0,01 \\
\hline
\centering Минимальная величина максимального из фазных токов (LIсим)  & \centering Iмакс> & \centering 0,01...2,00 & \centering о.е. & \centering \arraybackslash 0,01 \\
\hline
\centering Выдержка времени на срабатывание по критерию обрыва фазы (Tобрыв)  & \centering T1 & \centering 0,00...100,00 & \centering с & \centering \arraybackslash 0,01 \\
\hline
\centering Выдержка времени на срабатывание по асимметрии токов (Tасим)  & \centering T2 & \centering 0,00...100,00 & \centering с & \centering \arraybackslash 0,01 \\
\hline
%===t1
\end{longtable}
\normalsize


%%%%%%%%%%%%%%%%%%%%%%%%%%%%%%%%%%%%%%%%%%%%%%%%%%%%%%%%%% ЗПО %%%%%%%%%%%%%%%%%%%%%%%%%%%%%%%%%%%%%%%%%%%%%%%%%%%%%%%%%%%%%%%%%%
\par
В таблице \ref{app_set:zpo} приведены параметры, необходимые для настройки функции.

\footnotesize
\begin{longtable}{|>{\centering\arraybackslash}m{5.3cm}|>{\centering\arraybackslash}m{3.3cm}|>{\centering\arraybackslash}m{4.2cm}|>{\centering\arraybackslash}m{1.8cm}|>{\centering\arraybackslash}m{1cm}|}
\caption{Параметры для настройки ЗПО\hfill\vspace{-0.5\baselineskip}}\label{app_set:zpo}\\ 
\hline
\rowcolor{gray!20}
Параметр (Параметр на ИЧМ) & Условное обозначение на схеме & Значение/ Диапазон & Единица измерения & Шаг \\ 
\hline
\endfirsthead
%\caption{\emph{Продолжение\hfill\vspace{-0.5\baselineskip}}} \\ % Отступ обозначения таблицы от таблицы и выравнивание надписи слева
\caption*{\hspace{3pt}\emph{Продолжение таблицы \ref{app_set:zpo}\hfill\vspace{-0.5\baselineskip}}} \\ % сделано по ГОСТ 2.105 п.6.8.7
\hline
\rowcolor{gray!20}
Параметр (Параметр на ИЧМ) & Условное обозначение на схеме & Значение/ Диапазон & Единица измерения & Шаг \\ 
%\hline
\endhead
%\hline
\endfoot
%\hline
\endlastfoot
%==+t1*PALC1|EQPALC
\centering Ввод ЗПО в работу (Ввод\_функции)  & \centering SGF1 & \centering Не предусмотрено \\ Предусмотрено & \centering -- & \centering \arraybackslash -- \\
\hline
\centering Контроль тока (Контр\_тока)  & \centering SGF2 & \centering Не предусмотрено \\ Предусмотрено & \centering -- & \centering \arraybackslash -- \\
\hline
\centering Контроль температуры масла (Контр\_темп\_масла)  & \centering SGF3 & \centering Не предусмотрено \\ Предусмотрено & \centering -- & \centering \arraybackslash -- \\
\hline
\centering Выдержка времени срабатывания (Tср)  & \centering T1 & \centering 0...3600 & \centering с & \centering \arraybackslash 1 \\
\hline
%===t1
\end{longtable}
\normalsize

%%%%%%%%%%%%%%%%%%%%%%%%%%%%%%%%%%%%%%%%%%%%%%%%%%%%%%%%%% ТО ЗПО %%%%%%%%%%%%%%%%%%%%%%%%%%%%%%%%%%%%%%%%%%%%%%%%%%%%%%%%%%%%%%%%%%
\par
В таблице \ref{app_set:tozpo} приведены параметры, необходимые для настройки ТО ЗПО.

\footnotesize
\begin{longtable}{|>{\centering\arraybackslash}m{5.3cm}|>{\centering\arraybackslash}m{3.3cm}|>{\centering\arraybackslash}m{4.2cm}|>{\centering\arraybackslash}m{1.8cm}|>{\centering\arraybackslash}m{1cm}|}
\caption{Параметры для настройки функций в составе ТО ЗПО\hfill\vspace{-0.5\baselineskip}}\label{app_set:tozpo}\\ 
\hline
\rowcolor{gray!20}
Параметр (Параметр на ИЧМ) & Условное обозначение на схеме & Значение/ Диапазон & Единица измерения & Шаг \\ 
\hline
\endfirsthead
%\caption{\emph{Продолжение\hfill\vspace{-0.5\baselineskip}}} \\ % Отступ обозначения таблицы от таблицы и выравнивание надписи слева
\caption*{\hspace{3pt}\emph{Продолжение таблицы \ref{app_set:tozpo}\hfill\vspace{-0.5\baselineskip}}} \\ % сделано по ГОСТ 2.105 п.6.8.7
\hline
\rowcolor{gray!20}
Параметр (Параметр на ИЧМ) & Условное обозначение на схеме & Значение/ Диапазон & Единица измерения & Шаг \\ 
%\hline
\endhead
%\hline
\endfoot
%\hline
\endlastfoot
%==+t1*HVPTOC1|TPALC
\multicolumn{5}{|c|}{ ТО ЗПО ВН } \\ \hline 
\centering Ввод функции в работу (Ввод\_функции)  & \centering SGF1 & \centering Не предусмотрено \\ Предусмотрено & \centering -- & \centering \arraybackslash -- \\
\hline
\centering Ток срабатывания (Iср)  & \centering I> & \centering 0,10...30,00 & \centering о.е. & \centering \arraybackslash 0,01 \\
\hline
\multicolumn{5}{|c|}{ ТО ЗПО СН (НН1) } \\ \hline 
\centering Ввод функции в работу (Ввод\_функции)  & \centering SGF1 & \centering Не предусмотрено \\ Предусмотрено & \centering -- & \centering \arraybackslash -- \\
\hline
\centering Ток срабатывания (Iср)  & \centering I> & \centering 0,10...30,00 & \centering о.е. & \centering \arraybackslash 0,01 \\
\hline
\multicolumn{5}{|c|}{ ТО ЗПО НН (НН2) } \\ \hline 
\centering Ввод функции в работу (Ввод\_функции)  & \centering SGF1 & \centering Не предусмотрено \\ Предусмотрено & \centering -- & \centering \arraybackslash -- \\
\hline
\centering Ток срабатывания (Iср)  & \centering I> & \centering 0,10...30,00 & \centering о.е. & \centering \arraybackslash 0,01 \\
\hline
%===t1
\end{longtable}
\normalsize


%%%%%%%%%%%%%%%%%%%%%%%%%%%%%%%%%%%%%%%%%%%%%%%%%%%%%%%%%% ЗП %%%%%%%%%%%%%%%%%%%%%%%%%%%%%%%%%%%%%%%%%%%%%%%%%%%%%%%%%%%%%%%%%%
\par
В таблице \ref{app_set:zp} приведены параметры, необходимые для настройки ЗП.

\footnotesize
\begin{longtable}{|>{\centering\arraybackslash}m{5.3cm}|>{\centering\arraybackslash}m{3.3cm}|>{\centering\arraybackslash}m{4.2cm}|>{\centering\arraybackslash}m{1.8cm}|>{\centering\arraybackslash}m{1cm}|}
\caption{Параметры для настройки функций в составе ЗП\hfill\vspace{-0.5\baselineskip}}\label{app_set:zp}\\ 
\hline
\rowcolor{gray!20}
Параметр (Параметр на ИЧМ) & Условное обозначение на схеме & Значение/ Диапазон & Единица измерения & Шаг \\ 
\hline
\endfirsthead
%\caption{\emph{Продолжение\hfill\vspace{-0.5\baselineskip}}} \\ % Отступ обозначения таблицы от таблицы и выравнивание надписи слева
\caption*{\hspace{3pt}\emph{Продолжение таблицы \ref{app_set:zp}\hfill\vspace{-0.5\baselineskip}}} \\ % сделано по ГОСТ 2.105 п.6.8.7
\hline
\rowcolor{gray!20}
Параметр (Параметр на ИЧМ) & Условное обозначение на схеме & Значение/ Диапазон & Единица измерения & Шаг \\ 
%\hline
\endhead
%\hline
\endfoot
%\hline
\endlastfoot
%==+t1*HVPTOC1|TOVCTOC
\multicolumn{5}{|c|}{ ЗП ВН } \\ \hline 
\centering Ввод функции в работу (Ввод\_функции)  & \centering SGF1 & \centering Не предусмотрено \\ Предусмотрено & \centering -- & \centering \arraybackslash -- \\
\hline
\centering Ток срабатывания (Iср)  & \centering I> & \centering 0,10...5,00 & \centering о.е. & \centering \arraybackslash 0,01 \\
\hline
\centering Выдержка времени срабатывания (Tср)  & \centering T1 & \centering 0...3600000 & \centering мс & \centering \arraybackslash 100 \\
\hline
\multicolumn{5}{|c|}{ ЗП СН (НН1) } \\ \hline 
\centering Ввод функции в работу (Ввод\_функции)  & \centering SGF1 & \centering Не предусмотрено \\ Предусмотрено & \centering -- & \centering \arraybackslash -- \\
\hline
\centering Ток срабатывания (Iср)  & \centering I> & \centering 0,10...5,00 & \centering о.е. & \centering \arraybackslash 0,01 \\
\hline
\centering Выдержка времени срабатывания (Tср)  & \centering T1 & \centering 0...3600000 & \centering мс & \centering \arraybackslash 100 \\
\hline
\multicolumn{5}{|c|}{ ЗП НН (НН2) } \\ \hline 
\centering Ввод функции в работу (Ввод\_функции)  & \centering SGF1 & \centering Не предусмотрено \\ Предусмотрено & \centering -- & \centering \arraybackslash -- \\
\hline
\centering Ток срабатывания (Iср)  & \centering I> & \centering 0,10...5,00 & \centering о.е. & \centering \arraybackslash 0,01 \\
\hline
\centering Выдержка времени срабатывания (Tср)  & \centering T1 & \centering 0...3600000 & \centering мс & \centering \arraybackslash 100 \\
\hline
%===t1
\end{longtable}
\normalsize


%%%%%%%%%%%%%%%%%%%%%%%%%%%%%%%%%%%%%%%%%%%%%%%%%%%%%%%%%% РТПО %%%%%%%%%%%%%%%%%%%%%%%%%%%%%%%%%%%%%%%%%%%%%%%%%%%%%%%%%%%%%%%%%%
\par
В таблице \ref{rtpo:tbl1} приведены параметры, необходимые для настройки РТПО.

\footnotesize
\begin{longtable}{|>{\centering\arraybackslash}m{5.3cm}|>{\centering\arraybackslash}m{3.3cm}|>{\centering\arraybackslash}m{4.2cm}|>{\centering\arraybackslash}m{1.8cm}|>{\centering\arraybackslash}m{1cm}|}
\caption{Параметры для настройки функций в составе РТПО\hfill\vspace{-0.5\baselineskip}}\label{rtpo:tbl1}\\ 
\hline
\rowcolor{gray!20}
Параметр (Параметр на ИЧМ) & Условное обозначение на схеме & Значение/ Диапазон & Единица измерения & Шаг \\ 
\hline
\endfirsthead
%\caption{\emph{Продолжение\hfill\vspace{-0.5\baselineskip}}} \\ % Отступ обозначения таблицы от таблицы и выравнивание надписи слева
\caption*{\hspace{3pt}\emph{Продолжение таблицы \ref{rtpo:tbl1}\hfill\vspace{-0.5\baselineskip}}} \\ % сделано по ГОСТ 2.105 п.6.8.7
\hline
\rowcolor{gray!20}
Параметр (Параметр на ИЧМ) & Условное обозначение на схеме & Значение/ Диапазон & Единица измерения & Шаг \\ 
%\hline
\endhead
%\hline
\endfoot
%\hline
\endlastfoot
%==+t1*HVPTOC1|STRTPALC
\multicolumn{5}{|c|}{ РТПО ВН } \\ \hline 
\centering Ввод функции в работу (Ввод\_функции)  & \centering SGF1 & \centering Не предусмотрено \\ Предусмотрено & \centering -- & \centering \arraybackslash -- \\
\hline
\centering Ток срабатывания (Iср)  & \centering I> & \centering 0,10...30,00 & \centering о.е. & \centering \arraybackslash 0,01 \\
\hline
\multicolumn{5}{|c|}{ РТПО СН (НН1) } \\ \hline 
\centering Ввод функции в работу (Ввод\_функции)  & \centering SGF1 & \centering Не предусмотрено \\ Предусмотрено & \centering -- & \centering \arraybackslash -- \\
\hline
\centering Ток срабатывания (Iср)  & \centering I> & \centering 0,10...30,00 & \centering о.е. & \centering \arraybackslash 0,01 \\
\hline
\multicolumn{5}{|c|}{ РТПО НН (НН2) } \\ \hline 
\centering Ввод функции в работу (Ввод\_функции)  & \centering SGF1 & \centering Не предусмотрено \\ Предусмотрено & \centering -- & \centering \arraybackslash -- \\
\hline
\centering Ток срабатывания (Iср)  & \centering I> & \centering 0,10...30,00 & \centering о.е. & \centering \arraybackslash 0,01 \\
\hline
%===t1
\end{longtable}
\normalsize


%%%%%%%%%%%%%%%%%%%%%%%%%%%%%%%%%%%%%%%%%%%%%%%%%%%%%%%%%% ТК ЗДЗ %%%%%%%%%%%%%%%%%%%%%%%%%%%%%%%%%%%%%%%%%%%%%%%%%%%%%%%%%%%%%%%%%%

В таблице \ref{tkzdz:tbl1} приведены параметры, необходимые для настройки ТК ЗДЗ.

\footnotesize
\begin{longtable}{|>{\centering\arraybackslash}m{5.3cm}|>{\centering\arraybackslash}m{3.3cm}|>{\centering\arraybackslash}m{4.2cm}|>{\centering\arraybackslash}m{1.8cm}|>{\centering\arraybackslash}m{1cm}|}
\caption{Параметры для настройки ТК ЗДЗ\hfill\vspace{-0.5\baselineskip}}\label{tkzdz:tbl1}\\ 
\hline
\rowcolor{gray!20}
Параметр (Параметр на ИЧМ) & Условное обозначение на схеме & Значение/ Диапазон & Единица измерения & Шаг \\ 
\hline
\endfirsthead
%\caption{\emph{Продолжение\hfill\vspace{-0.5\baselineskip}}} \\ % Отступ обозначения таблицы от таблицы и выравнивание надписи слева
\caption*{\hspace{3pt}\emph{Продолжение таблицы \ref{tkzdz:tbl1}\hfill\vspace{-0.5\baselineskip}}} \\ % сделано по ГОСТ 2.105 п.6.8.7
\hline
\rowcolor{gray!20}
Параметр (Параметр на ИЧМ) & Условное обозначение на схеме & Значение/ Диапазон & Единица измерения & Шаг \\ 
%\hline
\endhead
%\hline
\endfoot
%\hline
\endlastfoot
%==+t1*PTOC1|LVARCTOC
\centering Ввод функции в работу (Ввод\_функции)  & \centering SGF1 & \centering Не предусмотрено \\ Предусмотрено & \centering -- & \centering \arraybackslash -- \\
\hline
\centering Ток срабатывания (Iср)  & \centering I> & \centering 0,10...30,00 & \centering о.е. & \centering \arraybackslash 0,01 \\
\hline
%===t1
\end{longtable}
\normalsize


%%%%%%%%%%%%%%%%%%%%%%%%%%%%%%%%%%%%%%%%%%%%%%%%%%%%%%%%%% ТО РПН %%%%%%%%%%%%%%%%%%%%%%%%%%%%%%%%%%%%%%%%%%%%%%%%%%%%%%%%%%%%%%%%%%

\par
В таблице \ref{torpn:tbl1} приведены параметры, необходимые для настройки ТО РПН.

\footnotesize
\begin{longtable}{|>{\centering\arraybackslash}m{5.3cm}|>{\centering\arraybackslash}m{3.3cm}|>{\centering\arraybackslash}m{4.2cm}|>{\centering\arraybackslash}m{1.8cm}|>{\centering\arraybackslash}m{1cm}|}
\caption{Параметры для настройки ТО РПН\hfill\vspace{-0.5\baselineskip}}\label{torpn:tbl1}\\ 
\hline
\rowcolor{gray!20}
Параметр (Параметр на ИЧМ) & Условное обозначение на схеме & Значение/ Диапазон & Единица измерения & Шаг \\ 
\hline
\endfirsthead
\caption*{\hspace{3pt}\emph{Продолжение таблицы \ref{torpn:tbl1}\hfill\vspace{-0.5\baselineskip}}} \\ % сделано по ГОСТ 2.105 п.6.8.7
\hline
\rowcolor{gray!20}
Параметр (Параметр на ИЧМ) & Условное обозначение на схеме & Значение/ Диапазон & Единица измерения & Шаг \\ 
%\hline
\endhead
%\hline
\endfoot
%\hline
\endlastfoot
%==+t1*PTOC1|LTCBLKTOC
\centering Ввод функции в работу (Ввод\_функции)  & \centering SGF1 & \centering Не предусмотрено \\ Предусмотрено & \centering -- & \centering \arraybackslash -- \\
\hline
\centering Ток срабатывания (Iср)  & \centering I> & \centering 0,10...30,00 & \centering о.е. & \centering \arraybackslash 0,01 \\
\hline
%===t1
\end{longtable}
\normalsize


%%%%%%%%%%%%%%%%%%%%%%%%%%%%%%%%%%%%%%%%%%%%%%%%%%%%%%%%%% ЛО ГЗоткл %%%%%%%%%%%%%%%%%%%%%%%%%%%%%%%%%%%%%%%%%%%%%%%%%%%%%%%%%%%%%%%%%%
\par

В таблице \ref{logzotkl:tbl1} приведены параметры, необходимые для настройки ЛО ГЗоткл.

\footnotesize
\begin{longtable}{|>{\centering\arraybackslash}m{5.3cm}|>{\centering\arraybackslash}m{3.3cm}|>{\centering\arraybackslash}m{4.2cm}|>{\centering\arraybackslash}m{1.8cm}|>{\centering\arraybackslash}m{1cm}|}
\caption{Параметры для настройки логики отключения при срабатывании отключающего контакта газового реле трансформатора\hfill\vspace{-0.5\baselineskip}}\label{logzotkl:tbl1}\\ 
\hline
\rowcolor{gray!20}
Параметр (Параметр на ИЧМ) & Условное обозначение на схеме & Значение/ Диапазон & Единица измерения & Шаг \\ 
\hline
\endfirsthead
\caption*{\hspace{3pt}\emph{Продолжение таблицы \ref{logzotkl:tbl1}\hfill\vspace{-0.5\baselineskip}}} \\ % сделано по ГОСТ 2.105 п.6.8.7
\hline
\rowcolor{gray!20}
Параметр (Параметр на ИЧМ) & Условное обозначение на схеме & Значение/ Диапазон & Единица измерения & Шаг \\ 
%\hline
\endhead
%\hline
\endfoot
%\hline
\endlastfoot
%==+t1*PTRC1|TTRGASLGC
\multicolumn{5}{|c|}{ БУ } \\ \hline 
\centering Выдержка времени срабатывания на блокировку (Тбл)  & \centering T1 & \centering 0...30000 & \centering мс & \centering \arraybackslash 10 \\
\hline
\multicolumn{5}{|c|}{ ЛО } \\ \hline 
\centering Ввод функции в работу (Ввод\_функции)  & \centering SGF1 & \centering Не предусмотрено \\ Предусмотрено & \centering -- & \centering \arraybackslash -- \\
\hline
\centering Блокировка от низкой изоляции (Блок\_от\_НИ)  & \centering SGF2 & \centering Не предусмотрено \\ Предусмотрено & \centering -- & \centering \arraybackslash -- \\
\hline
%===t1
\end{longtable}
\normalsize


%%%%%%%%%%%%%%%%%%%%%%%%%%%%%%%%%%%%%%%%%%%%%%%%%%%%%%%%%% ЛО ГЗсигн %%%%%%%%%%%%%%%%%%%%%%%%%%%%%%%%%%%%%%%%%%%%%%%%%%%%%%%%%%%%%%%%%%
\par

В таблице \ref{logzsign:tbl1} приведены параметры, необходимые для настройки ЛО ГЗсигн.

\footnotesize
\begin{longtable}{|>{\centering\arraybackslash}m{5.3cm}|>{\centering\arraybackslash}m{3.3cm}|>{\centering\arraybackslash}m{4.2cm}|>{\centering\arraybackslash}m{1.8cm}|>{\centering\arraybackslash}m{1cm}|}
\caption{Параметры для настройки логики отключения при срабатывании сигнального контакта газового реле трансформатора\hfill\vspace{-0.5\baselineskip}}\label{logzsign:tbl1}\\ 
\hline
\rowcolor{gray!20}
Параметр (Параметр на ИЧМ) & Условное обозначение на схеме & Значение/ Диапазон & Единица измерения & Шаг \\ 
\hline
\endfirsthead
\caption*{\hspace{3pt}\emph{Продолжение таблицы \ref{logzsign:tbl1}\hfill\vspace{-0.5\baselineskip}}} \\ % сделано по ГОСТ 2.105 п.6.8.7
\hline
\rowcolor{gray!20}
Параметр (Параметр на ИЧМ) & Условное обозначение на схеме & Значение/ Диапазон & Единица измерения & Шаг \\ 
%\hline
\endhead
%\hline
\endfoot
%\hline
\endlastfoot
%==+t1*PTRC1|TALMGASLGC
\multicolumn{5}{|c|}{ ЛО } \\ \hline 
\centering Ввод функции в работу (Ввод\_функции)  & \centering SGF1 & \centering Не предусмотрено \\ Предусмотрено & \centering -- & \centering \arraybackslash -- \\
\hline
\centering Выдержка времени срабатывания на блокировку (Тбл)  & \centering T1 & \centering 0...30000 & \centering мс & \centering \arraybackslash 10 \\
\hline
\centering Блокировка от низкой изоляции (Блок\_от\_НИ)  & \centering SGF2 & \centering Не предусмотрено \\ Предусмотрено & \centering -- & \centering \arraybackslash -- \\
\hline
%===t1
\end{longtable}
\normalsize

%%%%%%%%%%%%%%%%%%%%%%%%%%%%%%%%%%%%%%%%%%%%%%%%%%%%%%%%%% ЛО ГЗ РПН %%%%%%%%%%%%%%%%%%%%%%%%%%%%%%%%%%%%%%%%%%%%%%%%%%%%%%%%%%%%%%%%%%


\par

В таблице \ref{logzrpn:tbl1} приведены параметры, необходимые для настройки ЛО ГЗ РПН.

\footnotesize
\begin{longtable}{|>{\centering\arraybackslash}m{5.3cm}|>{\centering\arraybackslash}m{3.3cm}|>{\centering\arraybackslash}m{4.2cm}|>{\centering\arraybackslash}m{1.8cm}|>{\centering\arraybackslash}m{1cm}|}
\caption{Параметры для настройки логики отключения при срабатывании струйного реле РПН трансформатора\hfill\vspace{-0.5\baselineskip}}\label{logzrpn:tbl1}\\ 
\hline
\rowcolor{gray!20}
Параметр (Параметр на ИЧМ) & Условное обозначение на схеме & Значение/ Диапазон & Единица измерения & Шаг \\ 
\hline
\endfirsthead
\caption*{\hspace{3pt}\emph{Продолжение таблицы \ref{logzrpn:tbl1}\hfill\vspace{-0.5\baselineskip}}} \\ % сделано по ГОСТ 2.105 п.6.8.7
\hline
\rowcolor{gray!20}
Параметр (Параметр на ИЧМ) & Условное обозначение на схеме & Значение/ Диапазон & Единица измерения & Шаг \\ 
%\hline
\endhead
%\hline
\endfoot
%\hline
\endlastfoot
%==+t1*PTRC1|TLTCGASLGC
\multicolumn{5}{|c|}{ БУ } \\ \hline 
\centering Выдержка времени срабатывания на блокировку (Тбл)  & \centering T1 & \centering 0...30000 & \centering мс & \centering \arraybackslash 10 \\
\hline
\multicolumn{5}{|c|}{ ЛО } \\ \hline 
\centering Ввод функции в работу (Ввод\_функции)  & \centering SGF1 & \centering Не предусмотрено \\ Предусмотрено & \centering -- & \centering \arraybackslash -- \\
\hline
\centering Блокировка от низкой изоляции (Блок\_от\_НИ)  & \centering SGF2 & \centering Не предусмотрено \\ Предусмотрено & \centering -- & \centering \arraybackslash -- \\
\hline
%===t1
\end{longtable}
\normalsize



%%%%%%%%%%%%%%%%%%%%%%%%%%%%%%%%%%%%%%%%%%%%%%%%%%%%%%%%%% ЛО ТЗ %%%%%%%%%%%%%%%%%%%%%%%%%%%%%%%%%%%%%%%%%%%%%%%%%%%%%%%%%%%%%%%%%%
\par

В таблице \ref{lotz:tbl1} приведены параметры, необходимые для настройки ЛО ТЗ.
\footnotesize
\begin{longtable}{|>{\centering\arraybackslash}m{5.3cm}|>{\centering\arraybackslash}m{3.3cm}|>{\centering\arraybackslash}m{4.2cm}|>{\centering\arraybackslash}m{1.8cm}|>{\centering\arraybackslash}m{1cm}|}
\caption{Общие параметры для настройки ЛО ТЗ\hfill\vspace{-0.5\baselineskip}}\label{lotz:tbl1}\\ 
\hline
\rowcolor{gray!20}
Параметр (Параметр на ИЧМ) & Условное обозначение на схеме & Значение/ Диапазон & Единица измерения & Шаг \\ 
\hline
\endfirsthead
%\caption{\emph{Продолжение\hfill\vspace{-0.5\baselineskip}}} \\ % Отступ обозначения таблицы от таблицы и выравнивание надписи слева
\caption*{\hspace{3pt}\emph{Продолжение таблицы \ref{lotz:tbl1}\hfill\vspace{-0.5\baselineskip}}} \\ % сделано по ГОСТ 2.105 п.6.8.7
\hline
\rowcolor{gray!20}
Параметр (Параметр на ИЧМ) & Условное обозначение на схеме & Значение/ Диапазон & Единица измерения & Шаг \\ 
%\hline
\endhead
%\hline
\endfoot
%\hline
\endlastfoot
%==+t1*LLN0|APTTECHLGC
\multicolumn{5}{|c|}{ БУ } \\ \hline 
\centering Выдержка времени срабатывания на блокировку (Тбл)  & \centering T1 & \centering 0...30000 & \centering мс & \centering \arraybackslash 10 \\
\hline
\multicolumn{5}{|c|}{ ЛО ДТм } \\ \hline 
\centering Ввод функции в работу (Ввод\_функции)  & \centering SGF1 & \centering Не предусмотрено \\ Предусмотрено & \centering -- & \centering \arraybackslash -- \\
\hline
\centering Блокировка от низкой изоляции (Блок\_от\_НИ)  & \centering SGF2 & \centering Не предусмотрено \\ Предусмотрено & \centering -- & \centering \arraybackslash -- \\
\hline
\multicolumn{5}{|c|}{ ЛО ДТо } \\ \hline 
\centering Ввод функции в работу (Ввод\_функции)  & \centering SGF1 & \centering Не предусмотрено \\ Предусмотрено & \centering -- & \centering \arraybackslash -- \\
\hline
\centering Блокировка от низкой изоляции (Блок\_от\_НИ)  & \centering SGF2 & \centering Не предусмотрено \\ Предусмотрено & \centering -- & \centering \arraybackslash -- \\
\hline
\multicolumn{5}{|c|}{ ЛО РД } \\ \hline 
\centering Ввод функции в работу (Ввод\_функции)  & \centering SGF1 & \centering Не предусмотрено \\ Предусмотрено & \centering -- & \centering \arraybackslash -- \\
\hline
\centering Блокировка от низкой изоляции (Блок\_от\_НИ)  & \centering SGF2 & \centering Не предусмотрено \\ Предусмотрено & \centering -- & \centering \arraybackslash -- \\
\hline
%===t1
\end{longtable}
\normalsize


%%%%%%%%%%%%%%%%%%%%%%%%%%%%%%%%%%%%%%%%%%%%%%%%%%%%%%%%%% ЛО ТС %%%%%%%%%%%%%%%%%%%%%%%%%%%%%%%%%%%%%%%%%%%%%%%%%%%%%%%%%%%%%%%%%%
\par

В таблице \ref{lotspk:tbl1} приведены параметры, необходимые для настройки ЛО ТС.

\footnotesize
\begin{longtable}{|>{\centering\arraybackslash}m{5.3cm}|>{\centering\arraybackslash}m{3.3cm}|>{\centering\arraybackslash}m{4.2cm}|>{\centering\arraybackslash}m{1.8cm}|>{\centering\arraybackslash}m{1cm}|}
\caption{Параметры для настройки ЛО ПК\hfill\vspace{-0.5\baselineskip}}\label{lotspk:tbl1}\\ 
\hline
\rowcolor{gray!20}
Параметр (Параметр на ИЧМ) & Условное обозначение на схеме & Значение/ Диапазон & Единица измерения & Шаг \\ 
\hline
\endfirsthead
\caption*{\hspace{3pt}\emph{Продолжение таблицы \ref{lotspk:tbl1}\hfill\vspace{-0.5\baselineskip}}} \\ % сделано по ГОСТ 2.105 п.6.8.7
\hline
\rowcolor{gray!20}
Параметр (Параметр на ИЧМ) & Условное обозначение на схеме & Значение/ Диапазон & Единица измерения & Шаг \\ 
%\hline
\endhead
%\hline
\endfoot
%\hline
\endlastfoot
%==+t1*PRVLVPTRC1|ALMTECHLGC_UIRZ
\multicolumn{5}{|c|}{ ЛО ПК } \\ \hline 
\centering Ввод функции в работу (Ввод\_функции)  & \centering SGF1 & \centering Не предусмотрено \\ Предусмотрено & \centering -- & \centering \arraybackslash -- \\
\hline
\multicolumn{5}{|c|}{ ЛО ОК } \\ \hline 
\centering Ввод функции в работу (Ввод\_функции)  & \centering SGF1 & \centering Не предусмотрено \\ Предусмотрено & \centering -- & \centering \arraybackslash -- \\
\hline
\multicolumn{5}{|c|}{ ЛО ДУм расширителя } \\ \hline 
\centering Ввод функции в работу (Ввод\_функции)  & \centering SGF1 & \centering Не предусмотрено \\ Предусмотрено & \centering -- & \centering \arraybackslash -- \\
\hline
%===t1
\end{longtable}
\normalsize



%%%%%%%%%%%%%%%%%%%%%%%%%%%%%%%%%%%%%%%%%%%%%%%%%%%%%%%%%% ЛО Т %%%%%%%%%%%%%%%%%%%%%%%%%%%%%%%%%%%%%%%%%%%%%%%%%%%%%%%%%%%%%%%%%%

\par
В таблице \ref{loobsh:tbl1} приведены параметры, необходимые для настройки ЛО Т.

\footnotesize
\begin{longtable}{|>{\centering\arraybackslash}m{5.3cm}|>{\centering\arraybackslash}m{3.3cm}|>{\centering\arraybackslash}m{4.2cm}|>{\centering\arraybackslash}m{1.8cm}|>{\centering\arraybackslash}m{1cm}|}
\caption{Параметры для настройки ЛО Т\hfill\vspace{-0.5\baselineskip}}\label{loobsh:tbl1}\\ 
\hline
\rowcolor{gray!20}
Параметр (Параметр на ИЧМ) & Условное обозначение на схеме & Значение/ Диапазон & Единица измерения & Шаг \\ 
\hline
\endfirsthead
%\caption{\emph{Продолжение\hfill\vspace{-0.5\baselineskip}}} \\ % Отступ обозначения таблицы от таблицы и выравнивание надписи слева
\caption*{\hspace{3pt}\emph{Продолжение таблицы \ref{loobsh:tbl1}\hfill\vspace{-0.5\baselineskip}}} \\ % сделано по ГОСТ 2.105 п.6.8.7
\hline
\rowcolor{gray!20}
Параметр (Параметр на ИЧМ) & Условное обозначение на схеме & Значение/ Диапазон & Единица измерения & Шаг \\ 
%\hline
\endhead
%\hline
\endfoot
%\hline
\endlastfoot
%==+t1*PTRC1|TPRMOFFLVLGC
\multicolumn{5}{|c|}{ ЛО } \\ \hline 
\centering Ввод функции в работу (Ввод\_функции)  & \centering SGF1 & \centering Не предусмотрено \\ Предусмотрено & \centering -- & \centering \arraybackslash -- \\
\hline
\multicolumn{5}{|c|}{ ЗАПВ } \\ \hline 
\centering Ввод функции в работу (Ввод\_функции)  & \centering SGF1 & \centering Не предусмотрено \\ Предусмотрено & \centering -- & \centering \arraybackslash -- \\
\hline
%===t1
\end{longtable}
\normalsize



%%%%%%%%%%%%%%%%%%%%%%%%%%%%%%%%%%%%%%%%%%%%%%%%%%%%%%%%%% ЛО ВН %%%%%%%%%%%%%%%%%%%%%%%%%%%%%%%%%%%%%%%%%%%%%%%%%%%%%%%%%%%%%%%%%%

\par
В таблице \ref{lovn:tbl1} приведены параметры, необходимые для настройки ЛО ВН.

\footnotesize
\begin{longtable}{|>{\centering\arraybackslash}m{5.3cm}|>{\centering\arraybackslash}m{3.3cm}|>{\centering\arraybackslash}m{4.2cm}|>{\centering\arraybackslash}m{1.8cm}|>{\centering\arraybackslash}m{1cm}|}
\caption{Параметры для настройки ЛО ВН\hfill\vspace{-0.5\baselineskip}}\label{lovn:tbl1}\\ 
\hline
\rowcolor{gray!20}
Параметр (Параметр на ИЧМ) & Условное обозначение на схеме & Значение/ Диапазон & Единица измерения & Шаг \\ 
\hline
\endfirsthead
\caption*{\hspace{3pt}\emph{Продолжение таблицы \ref{lovn:tbl1}\hfill\vspace{-0.5\baselineskip}}} \\ % сделано по ГОСТ 2.105 п.6.8.7
\hline
\rowcolor{gray!20}
Параметр (Параметр на ИЧМ) & Условное обозначение на схеме & Значение/ Диапазон & Единица измерения & Шаг \\ 
%\hline
\endhead
%\hline
\endfoot
%\hline
\endlastfoot
%==+t1*HVCBPTRC1|HVTCBOFF
\multicolumn{5}{|c|}{ ЛО } \\ \hline 
\centering Ввод функции в работу (Ввод\_функции)  & \centering SGF1 & \centering Не предусмотрено \\ Предусмотрено & \centering -- & \centering \arraybackslash -- \\
\hline
\centering Длительность импульса (Тимп)  & \centering T1 & \centering 50...60000 & \centering мс & \centering \arraybackslash 10 \\
\hline
%===t1
\end{longtable}
\normalsize


%%%%%%%%%%%%%%%%%%%%%%%%%%%%%%%%%%%%%%%%%%%%%%%%%%%%%%%%%% ЛО НН %%%%%%%%%%%%%%%%%%%%%%%%%%%%%%%%%%%%%%%%%%%%%%%%%%%%%%%%%%%%%%%%%%

\par
В таблице \ref{lonn:tbl1} приведены параметры, необходимые для настройки ЛО НН.

\footnotesize
\begin{longtable}{|>{\centering\arraybackslash}m{5.3cm}|>{\centering\arraybackslash}m{3.3cm}|>{\centering\arraybackslash}m{4.2cm}|>{\centering\arraybackslash}m{1.8cm}|>{\centering\arraybackslash}m{1cm}|}
\caption{Параметры для настройки логики отключения выключателя\hfill\vspace{-0.5\baselineskip}}\label{lonn:tbl1}\\ 
\hline
\rowcolor{gray!20}
Параметр (Параметр на ИЧМ) & Условное обозначение на схеме & Значение/ Диапазон & Единица измерения & Шаг \\ 
\hline
\endfirsthead
%\caption{\emph{Продолжение\hfill\vspace{-0.5\baselineskip}}} \\ % Отступ обозначения таблицы от таблицы и выравнивание надписи слева
\caption*{\hspace{3pt}\emph{Продолжение таблицы \ref{lonn:tbl1}\hfill\vspace{-0.5\baselineskip}}} \\ % сделано по ГОСТ 2.105 п.6.8.7
\hline
\rowcolor{gray!20}
Параметр (Параметр на ИЧМ) & Условное обозначение на схеме & Значение/ Диапазон & Единица измерения & Шаг \\ 
%\hline
\endhead
%\hline
\endfoot
%\hline
\endlastfoot
%==+t1*LVCBPTRC1|LVTPRMCBOFF1
\multicolumn{5}{|c|}{ ЛО } \\ \hline 
\centering Ввод функции в работу (Ввод\_функции)  & \centering SGF1 & \centering Не предусмотрено \\ Предусмотрено & \centering -- & \centering \arraybackslash -- \\
\hline
\centering Длительность импульса (Тимп)  & \centering T1 & \centering 50...60000 & \centering мс & \centering \arraybackslash 10 \\
\hline
\multicolumn{5}{|c|}{ ЗАПВ } \\ \hline 
\centering Ввод функции в работу (Ввод\_функции)  & \centering SGF1 & \centering Не предусмотрено \\ Предусмотрено & \centering -- & \centering \arraybackslash -- \\
\hline
%===t1
\end{longtable}
\normalsize


%%%%%%%%%%%%%%%%%%%%%%%%%%%%%%%%%%%%%%%%%%%%%%%%%%%%%%%%%% УРОВ %%%%%%%%%%%%%%%%%%%%%%%%%%%%%%%%%%%%%%%%%%%%%%%%%%%%%%%%%%%%%%%%%%

\par
В таблице \ref{urov35:tbl1} приведены параметры, необходимые для настройки УРОВ.

\footnotesize
\begin{longtable}{|>{\centering\arraybackslash}m{5.3cm}|>{\centering\arraybackslash}m{3.3cm}|>{\centering\arraybackslash}m{4.2cm}|>{\centering\arraybackslash}m{1.8cm}|>{\centering\arraybackslash}m{1cm}|}
\caption{Параметры для настройки УРОВ\hfill\vspace{-0.5\baselineskip}}\label{urov35:tbl1}\\ 
\hline
\rowcolor{gray!20}
Параметр (Параметр на ИЧМ) & Условное обозначение на схеме & Значение/ Диапазон & Единица измерения & Шаг \\ 
\hline
\endfirsthead
\caption{\emph{Продолжение\hfill\vspace{-0.5\baselineskip}}} \\ % Отступ обозначения таблицы от таблицы и выравнивание надписи слева
\hline
\rowcolor{gray!20}
Параметр (Параметр на ИЧМ) & Условное обозначение на схеме & Значение/ Диапазон & Единица измерения & Шаг \\ 
%\hline
\endhead
%\hline
\endfoot
%\hline
\endlastfoot
%==+t1*GENRBRF1|DZT2_TPBRF
\multicolumn{5}{|c|}{ УРОВ } \\ \hline 
\centering Ввод функции в работу (Ввод\_функции)  & \centering SGF1 & \centering Не предусмотрено \\ Предусмотрено & \centering -- & \centering \arraybackslash -- \\
\hline
\centering Ток срабатывания (Iср)  & \centering I> & \centering 0,05...0,50 & \centering о.е. & \centering \arraybackslash 0,01 \\
\hline
\centering УРОВ с подхватом по току (Подхват\_по\_току)  & \centering SGF2 & \centering Не предусмотрено \\ Предусмотрено & \centering -- & \centering \arraybackslash -- \\
\hline
\centering Контроль по току при действии "на себя" (Контр\_тока\_на\_себя)  & \centering SGF4 & \centering Не предусмотрено \\ Предусмотрено & \centering -- & \centering \arraybackslash -- \\
\hline
\centering Выдержка времени срабатывания (Tср)  & \centering T1 & \centering 0...1000 & \centering мс & \centering \arraybackslash 10 \\
\hline
\centering Действие внешнего УРОВ на вышестоящий выключатель (Действ\_на\_выш\_выкл)  & \centering SGF3 & \centering Не предусмотрено \\ Предусмотрено & \centering -- & \centering \arraybackslash -- \\
\hline
\centering Контроль ЭМО (Контроль\_ЭМО)  & \centering SGF5 & \centering Не предусмотрено \\ Предусмотрено по ЭМО1 \\ Предусмотрено по ЭМО1 и ЭМО2 & \centering -- & \centering \arraybackslash -- \\
\hline
%===t1
\end{longtable}
\normalsize


%%%%%%%%%%%%%%%%%%%%%%%%%%%%%%%%%%%%%%%%%%%%%%%%%%%%%%%%%% ПС %%%%%%%%%%%%%%%%%%%%%%%%%%%%%%%%%%%%%%%%%%%%%%%%%%%%%%%%%%%%%%%%%%

\par
В таблице \ref{ps:tbl1} приведены параметры, необходимые для настройки ПС. 

\footnotesize
\begin{longtable}{|>{\centering\arraybackslash}m{5.3cm}|>{\centering\arraybackslash}m{3.3cm}|>{\centering\arraybackslash}m{4.2cm}|>{\centering\arraybackslash}m{1.8cm}|>{\centering\arraybackslash}m{1cm}|}
\caption{Параметры для настройки предупредительной сигнализации\hfill\vspace{-0.5\baselineskip}}\label{ps:tbl1}\\ 
\hline
\rowcolor{gray!20}
Параметр (Параметр на ИЧМ) & Условное обозначение на схеме & Значение/ Диапазон & Единица измерения & Шаг \\ 
\hline
\endfirsthead
%\caption{\emph{Продолжение\hfill\vspace{-0.5\baselineskip}}} \\ % Отступ обозначения таблицы от таблицы и выравнивание надписи слева
\caption*{\hspace{3pt}\emph{Продолжение таблицы \ref{ps:tbl1}\hfill\vspace{-0.5\baselineskip}}} \\ % сделано по ГОСТ 2.105 п.6.8.7
\hline
\rowcolor{gray!20}
Параметр (Параметр на ИЧМ) & Условное обозначение на схеме & Значение/ Диапазон & Единица измерения & Шаг \\ 
%\hline
\endhead
%\hline
\endfoot
%\hline
\endlastfoot
%==+t1*CALH1|DZT2_LVALH
\centering Контроль сигнала «ГЗ сигн» (Контр\_ГЗ\_сигн)  & \centering SGF1 & \centering Не предусмотрено \\ Предусмотрено & \centering -- & \centering \arraybackslash -- \\
\hline
\centering Контроль сигнала «Низ. изол. ГЗ» (Контр\_Низ\_изол\_ГЗ)  & \centering SGF2 & \centering Не предусмотрено \\ Предусмотрено & \centering -- & \centering \arraybackslash -- \\
\hline
\centering Контроль сигнала «ГЗ заблокирована» (Контр\_ГЗ\_блок)  & \centering SGF3 & \centering Не предусмотрено \\ Предусмотрено & \centering -- & \centering \arraybackslash -- \\
\hline
\centering Контроль сигнала «ТЗ сигн» (Контр\_ТЗ\_сигн)  & \centering SGF4 & \centering Не предусмотрено \\ Предусмотрено & \centering -- & \centering \arraybackslash -- \\
\hline
\centering Контроль сигнала «Низ. изол. ТЗ» (Контр\_Низ\_изол\_ТЗ)  & \centering SGF5 & \centering Не предусмотрено \\ Предусмотрено & \centering -- & \centering \arraybackslash -- \\
\hline
\centering Контроль сигнала «ТЗ заблокирована» (Контр\_ТЗ\_блок)  & \centering SGF6 & \centering Не предусмотрено \\ Предусмотрено & \centering -- & \centering \arraybackslash -- \\
\hline
\centering Контроль сигнала «ТС сигн» (Контр\_ТС\_сигн)  & \centering SGF7 & \centering Не предусмотрено \\ Предусмотрено & \centering -- & \centering \arraybackslash -- \\
\hline
\centering Контроль сигнала «ОТ сигн» (Контр\_ОТ\_сигн)  & \centering SGF8 & \centering Не предусмотрено \\ Предусмотрено & \centering -- & \centering \arraybackslash -- \\
\hline
\centering Контроль сигнала «ОТ НН сигн» (Контр\_ОТ\_НН\_сигн)  & \centering SGF9 & \centering Не предусмотрено \\ Предусмотрено & \centering -- & \centering \arraybackslash -- \\
\hline
\centering Контроль сигнала «Внеш. откл» (Контр\_Внеш\_откл)  & \centering SGF10 & \centering Не предусмотрено \\ Предусмотрено & \centering -- & \centering \arraybackslash -- \\
\hline
\centering Контроль сигнала «Вых. цепи разобраны» (Контр\_Вых\_цеп\_разоб)  & \centering SGF11 & \centering Не предусмотрено \\ Предусмотрено & \centering -- & \centering \arraybackslash -- \\
\hline
\centering Контроль сигнала «БИ выведены» (Контр\_БИ\_выведены)  & \centering SGF12 & \centering Не предусмотрено \\ Предусмотрено & \centering -- & \centering \arraybackslash -- \\
\hline
\centering Контроль общего внешнего сигнала (Контр\_общ\_внеш\_сигн)  & \centering SGF13 & \centering Не предусмотрено \\ Предусмотрено & \centering -- & \centering \arraybackslash -- \\
\hline
%===t1
\end{longtable}
\normalsize


%%%%%%%%%%%%%%%%%%%%%%%%%%%%%%%%%%%%%%%%%%%%%%%%%%%%%%%%%% СС %%%%%%%%%%%%%%%%%%%%%%%%%%%%%%%%%%%%%%%%%%%%%%%%%%%%%%%%%%%%%%%%%%

\par
В таблице \ref{ss:tbl1} приведены параметры, необходимые для настройки СС.

\footnotesize
\begin{longtable}{|>{\centering\arraybackslash}m{5.3cm}|>{\centering\arraybackslash}m{3.3cm}|>{\centering\arraybackslash}m{4.2cm}|>{\centering\arraybackslash}m{1.8cm}|>{\centering\arraybackslash}m{1cm}|}
\caption{Параметры для настройки функции сборки сигналов\hfill\vspace{-0.5\baselineskip}}\label{ss:tbl1}\\ 
\hline
\rowcolor{gray!20}
Параметр (Параметр на ИЧМ) & Условное обозначение на схеме & Значение/ Диапазон & Единица измерения & Шаг \\ 
\hline
\endfirsthead
%\caption{\emph{Продолжение\hfill\vspace{-0.5\baselineskip}}} \\ % Отступ обозначения таблицы от таблицы и выравнивание надписи слева
\caption*{\hspace{3pt}\emph{Продолжение таблицы \ref{ss:tbl1}\hfill\vspace{-0.5\baselineskip}}} \\ % сделано по ГОСТ 2.105 п.6.8.7
\hline
\rowcolor{gray!20}
Параметр (Параметр на ИЧМ) & Условное обозначение на схеме & Значение/ Диапазон & Единица измерения & Шаг \\ 
%\hline
\endhead
%\hline
\endfoot
%\hline
\endlastfoot
%==+t1*GAPC1|DZT2_SignAssembly
\multicolumn{5}{|c|}{ CC } \\ \hline 
\centering Контроль положения SA1 (Контр\_полож\_SA1)  & \centering SGF1 & \centering Не предусмотрено \\ Предусмотрено & \centering -- & \centering \arraybackslash -- \\
\hline
\centering Контроль положения SA2 (Контр\_полож\_SA2)  & \centering SGF2 & \centering Не предусмотрено \\ Предусмотрено & \centering -- & \centering \arraybackslash -- \\
\hline
\centering Контроль положения SA3 (Контр\_полож\_SA3)  & \centering SGF3 & \centering Не предусмотрено \\ Предусмотрено & \centering -- & \centering \arraybackslash -- \\
\hline
\centering Контроль положения SA4 (Контр\_полож\_SA4)  & \centering SGF4 & \centering Не предусмотрено \\ Предусмотрено & \centering -- & \centering \arraybackslash -- \\
\hline
\centering Контроль положения SG1 (Контр\_полож\_SG1)  & \centering SGF5 & \centering Не предусмотрено \\ Предусмотрено & \centering -- & \centering \arraybackslash -- \\
\hline
\centering Контроль положения SG2 (Контр\_полож\_SG2)  & \centering SGF6 & \centering Не предусмотрено \\ Предусмотрено & \centering -- & \centering \arraybackslash -- \\
\hline
\centering Контроль положения SG3 (Контр\_полож\_SG3)  & \centering SGF7 & \centering Не предусмотрено \\ Предусмотрено & \centering -- & \centering \arraybackslash -- \\
\hline
\centering Контроль опер. тока цепей ГЗ (Контр\_ОТ\_ГЗ)  & \centering SGF8 & \centering Не предусмотрено \\ Предусмотрено & \centering -- & \centering \arraybackslash -- \\
\hline
\centering Контроль опер. тока цепей ТЗ, ТС (Контр\_ОТ\_ТЗ\_ТС)  & \centering SGF9 & \centering Не предусмотрено \\ Предусмотрено & \centering -- & \centering \arraybackslash -- \\
\hline
\centering Контроль опер. тока цепей ЗДЗ НН1 (Контр\_ОТ\_ЗДЗ1)  & \centering SGF10 & \centering Не предусмотрено \\ Предусмотрено & \centering -- & \centering \arraybackslash -- \\
\hline
\centering Контроль опер. тока цепей ЗДЗ НН2 (Контр\_ОТ\_ЗДЗ2)  & \centering SGF11 & \centering Не предусмотрено \\ Предусмотрено & \centering -- & \centering \arraybackslash -- \\
\hline
\centering Контроль опер. тока цепей УРОВ НН1 (Контр\_ОТ\_УРОВ1)  & \centering SGF12 & \centering Не предусмотрено \\ Предусмотрено & \centering -- & \centering \arraybackslash -- \\
\hline
\centering Контроль опер. тока цепей УРОВ НН2 (Контр\_ОТ\_УРОВ2)  & \centering SGF13 & \centering Не предусмотрено \\ Предусмотрено & \centering -- & \centering \arraybackslash -- \\
\hline
%===t1
\end{longtable}
\normalsize
